次に、オブソレッセンス状態にあると判定された高被引用論文について、その群内で特徴差があるのかを確認するため階層クラスタリングを行った。クラスタリングに用いた特徴量は、a. 経過年に関する指標としてa-1:出版からピークまでの期間、a-2:ピークから被引用の十分な減衰が見られるまでの期間、b.被引用の寄与に関する指標としてb-1:ピーク年の被引用数、b-2:減衰年までの被引用数合計、c.被引用の安定性に関する指標としてc-1:毎年の被引用数の分散、c-2:毎年の被引用速度の分散、c-3:毎年の引用加速度の分散（以上、7指標）である。距離にはユークリッド距離を用い、各指標値にはノーマライズを行った。分散の計算においては先頭に0-パディングを行った5年平均を用いた。この結果を図\ref{fig:gr8}に示す。クラスタリングではノードが集中的に存在する領域といくつかのノードがまばらに見られる周辺領域が観察されるが、明確な分割構造は確認できなかった。
